\section{Formal Proofs of Core Theoretical Results}
\label{app:proofs}

\subsection{Proof of Theorem 1 (Zero-Shot Spectral Optimality of Res-SVD)}
\begin{proof}
Let dense weight matrix $W \in \mathbb{R}^{d_{\text{out}} \times d_{\text{in}}}$ have singular value decomposition $W = U \Sigma V^\top = \sum_{i=1}^d \sigma_i u_i v_i^\top$ with ordered singular values $\sigma_1 \ge \sigma_2 \ge \dots \ge \sigma_d \ge 0$.
The truncated base rank-$r$ SVD approximation is $W_{\text{base}} = \sum_{i=1}^r \sigma_i u_i v_i^\top$. By the Eckart-Young-Mirsky theorem, the Frobenius reconstruction error of the base model is:
\begin{equation}
\|W - W_{\text{base}}\|_F^2 = \left\| \sum_{i=r+1}^d \sigma_i u_i v_i^\top \right\|_F^2 = \sum_{i=r+1}^d \sigma_i^2.
\end{equation}
In standard LoRA, the adapter matrices are initialized such that $\gamma B A = 0$, yielding initial error $\|W - (W_{\text{base}} + 0)\|_F^2 = \sum_{i=r+1}^d \sigma_i^2$.

In Res-SVD, adapter matrices $A \in \mathbb{R}^{k \times d_{\text{in}}}$ and $B \in \mathbb{R}^{d_{\text{out}} \times k}$ are initialized directly from the next $k$ singular vectors:
\begin{equation}
B = U_{\text{tail}} \sqrt{\frac{\Sigma_{\text{tail}}}{\gamma}}, \quad A = \sqrt{\frac{\Sigma_{\text{tail}}}{\gamma}} V_{\text{tail}}^\top,
\end{equation}
where $U_{\text{tail}} = [u_{r+1}, \dots, u_{r+k}]$, $\Sigma_{\text{tail}} = \operatorname{diag}(\sigma_{r+1}, \dots, \sigma_{r+k})$, and $V_{\text{tail}} = [v_{r+1}, \dots, v_{r+k}]$.
The product of the adapters is:
\begin{equation}
\gamma B A = \gamma \left( U_{\text{tail}} \sqrt{\frac{\Sigma_{\text{tail}}}{\gamma}} \right) \left( \sqrt{\frac{\Sigma_{\text{tail}}}{\gamma}} V_{\text{tail}}^\top \right) = \sum_{i=r+1}^{r+k} \sigma_i u_i v_i^\top.
\end{equation}
The composite weight reconstruction $W_{\text{approx}} = W_{\text{base}} + \gamma B A$ is therefore:
\begin{equation}
W_{\text{approx}} = \sum_{i=1}^r \sigma_i u_i v_i^\top + \sum_{i=r+1}^{r+k} \sigma_i u_i v_i^\top = \sum_{i=1}^{r+k} \sigma_i u_i v_i^\top.
\end{equation}
The residual reconstruction error is:
\begin{equation}
\|W - W_{\text{approx}}\|_F^2 = \left\| \sum_{i=r+k+1}^d \sigma_i u_i v_i^\top \right\|_F^2 = \sum_{i=r+k+1}^d \sigma_i^2.
\end{equation}
Since $\sigma_i > 0$ for non-zero singular values, $\sum_{i=r+k+1}^d \sigma_i^2 < \sum_{i=r+1}^d \sigma_i^2$, establishing strict zero-shot reconstruction error minimization.
\end{proof}

\subsection{Proof of Theorem 2 (Steady-State Associative Memory Boundedness)}
\begin{proof}
Consider the recurrent linear attention state update equation:
\begin{equation}
S_t = \lambda_t S_{t-1} + \phi(K_t)^\top V_t, \quad S_0 = 0.
\end{equation}
Unrolling the recurrence from $t=1$ to $N$ yields:
\begin{equation}
S_N = \sum_{j=1}^N \left( \prod_{m=j+1}^N \lambda_m \right) \phi(K_j)^\top V_j,
\end{equation}
where $\prod_{m=N+1}^N \lambda_m \equiv 1$.
Applying the triangle inequality and sub-multiplicative property of the Frobenius matrix norm:
\begin{align}
\|S_N\|_F &\le \sum_{j=1}^N \left( \prod_{m=j+1}^N |\lambda_m| \right) \|\phi(K_j)^\top V_j\|_F \\
&\le \sum_{j=1}^N \left( \prod_{m=j+1}^N \lambda_m \right) \|\phi(K_j)\|_F \|V_j\|_F.
\end{align}
By assumption, inputs are bounded such that $\|\phi(K_j)\|_F \le M_K$ and $\|V_j\|_F \le M_V$. Furthermore, the decay rate satisfies $\lambda_t \le \lambda_{\max} < 1$ for all $t$. Thus:
\begin{equation}
\prod_{m=j+1}^N \lambda_m \le \lambda_{\max}^{N-j}.
\end{equation}
Substituting this upper bound:
\begin{equation}
\|S_N\|_F \le M_K M_V \sum_{j=1}^N \lambda_{\max}^{N-j} = M_K M_V \sum_{k=0}^{N-1} \lambda_{\max}^k.
\end{equation}
Evaluating the geometric series as $N \to \infty$:
\begin{equation}
\lim_{N \to \infty} \|S_N\|_F \le M_K M_V \sum_{k=0}^\infty \lambda_{\max}^k = \frac{M_K M_V}{1 - \lambda_{\max}} < \infty.
\end{equation}
Hence, the state Frobenius norm remains strictly bounded for sequences of arbitrary length.
\end{proof}

\subsection{Derivation of Fast Nystr\"om 2-Wasserstein Metric Equivalence}
\begin{proposition}
Let $\Sigma_1 = U_1 \Lambda_1 U_1^\top$ and $\Sigma_2 = U_2 \Lambda_2 U_2^\top$ be low-rank symmetric positive semi-definite covariance matrices with $U_i \in \operatorname{St}(r, d)$ and diagonal matrices $\Lambda_i \in \mathbb{R}^{r \times r}_{>0}$. Then:
\begin{equation}
\operatorname{Tr}\left(\left(\Sigma_1^{1/2} \Sigma_2 \Sigma_1^{1/2}\right)^{1/2}\right) = \left\| \Lambda_1^{1/2} U_1^\top U_2 \Lambda_2^{1/2} \right\|_*,
\end{equation}
where $\|\cdot\|_*$ denotes the matrix nuclear norm (sum of singular values).
\end{proposition}
\begin{proof}
Let $M = \Sigma_1^{1/2} \Sigma_2 \Sigma_1^{1/2}$. Since $\Sigma_i^{1/2} = U_i \Lambda_i^{1/2} U_i^\top$, we have:
\begin{equation}
M = U_1 \Lambda_1^{1/2} U_1^\top U_2 \Lambda_2 U_2^\top U_1 \Lambda_1^{1/2} U_1^\top.
\end{equation}
Define the $r \times r$ core cross-correlation matrix:
\begin{equation}
K = \Lambda_1^{1/2} U_1^\top U_2 \Lambda_2^{1/2} \in \mathbb{R}^{r \times r}.
\end{equation}
Notice that $K K^\top = \Lambda_1^{1/2} U_1^\top U_2 \Lambda_2 U_2^\top U_1 \Lambda_1^{1/2}$.
Therefore, $M = U_1 (K K^\top) U_1^\top$.
The non-zero eigenvalues of $M$ coincide exactly with the non-zero eigenvalues of $K K^\top$. Let the singular value decomposition of $K$ be $K = P S Q^\top$ where $S = \operatorname{diag}(s_1, \dots, s_r)$.
Then $K K^\top = P S^2 P^\top$, which implies:
\begin{equation}
M = (U_1 P) S^2 (U_1 P)^\top.
\end{equation}
Since $U_1 \in \operatorname{St}(r, d)$ and $P \in \mathcal{O}(r)$, the matrix $\tilde{U} = U_1 P$ satisfies $\tilde{U}^\top \tilde{U} = P^\top U_1^\top U_1 P = P^\top I_r P = I_r$, which is an orthonormal basis.
The matrix square root is $M^{1/2} = \tilde{U} S \tilde{U}^\top$.
Taking the trace:
\begin{equation}
\operatorname{Tr}(M^{1/2}) = \operatorname{Tr}(\tilde{U} S \tilde{U}^\top) = \operatorname{Tr}(S \tilde{U}^\top \tilde{U}) = \operatorname{Tr}(S) = \sum_{i=1}^r s_i.
\end{equation}
By definition, the sum of the singular values of $K$ is the nuclear norm $\|K\|_*$. Thus:
\begin{equation}
\operatorname{Tr}\left(\left(\Sigma_1^{1/2} \Sigma_2 \Sigma_1^{1/2}\right)^{1/2}\right) = \left\| \Lambda_1^{1/2} U_1^\top U_2 \Lambda_2^{1/2} \right\|_*,
\end{equation}
reducing the computational complexity from $\mathcal{O}(d^3)$ full matrix square roots to an $\mathcal{O}(r^2 d + r^3)$ SVD on the $r \times r$ core matrix $K$.
\end{proof}
